\chapter{Scope and Qualifications}

\section{Companion role}

This publication is EL-MW, the cumulative memory companion to Triptych's
print-first Ecclesiastical Latin curriculum.  It renders the same stable
lexeme, grammar-fact, recurrence, reconstruction, and composition registries
used by the individual packets.  It therefore supports packet teaching and
connected reading without replacing either one.

The principal destination is direct reading of the 1962 \emph{Missale
Romanum}.  The curriculum deliberately begins with ordinary project Latin and
extends through broader narrative, classical, biblical, patristic,
ecclesiastical, juridical, and poetic work, because durable Missal reading
depends on a real command of Latin rather than memorized liturgical fragments.
Composition develops from the first transparent clauses through controlled
imitation and independent prose.  Poetry is introduced in the common trunk
and becomes its own advanced branch.  No pronunciation system is taught here.

\section{Evidence, review, and rights}

The lexeme registry's concise English senses, pedagogical strands, recurrence
allocations, and every workbook prompt and model are project-created
instructional content.  They are not represented as quotations or as verified
authorial usage.  Received Latin belongs in the course's separately audited
passage registry with an identifiable work, edition, locus, and verification
state.  A full dictionary and grammar remain useful reference works, but the
forms, senses, construction cues, answer models, and recurrence sheets needed
to complete this workbook are printed here.

This is a language-study aid, not an official liturgical book, critical
edition, pronunciation course, or authority for a liturgical, theological,
textual, or canonical judgment.
No independent Latinist, liturgist, pedagogue, theologian, or ecclesiastical
reviewer is claimed.  The rights and distribution status of received works are
not altered by their citation as reference corpora.

